\centerline{\bf \Huge Числовые ребусы I}

\begin{flushright}
        \scriptsize{Добавлю вашу цитату}
\end{flushright}

\textbf{В ребусах с участием букв действует следующее правило: }

\textit{Одинаковым буквам должны соответствовать одинаковые цифры, разным буквам — разные цифры.}

\textbf{Решить ребус значит найти ВСЕ решения и доказать, что других нет!}

\problem{\textbf{1}} ББ + ББ = АБВ

\problem{\textbf{2}} ДА + ДА + ДА = ЕДА

\problem{\textbf{3}} УДАР + УДАР = ДРАКА

\problem{\textbf{4}} АИСТ + АИСТ + АИСТ + АИСТ = СТАЯ

\problem{\textbf{5}} ТОРГ $\times$ Г = ГРОТ

\problem{\textbf{6}} С + У + М + М + А = 11. Среди букв нет нуля. Чему может быть равна буква М?

\problem{\textbf{7}} ТЬМА + ЖУТЬ = УЖАС (6 решений)